% ID: SEL-D0063-OB-E01
% ESTADO: validada. S063.
\textbf{a)} Se sombrea en el primer cuadrante por debajo de $4x/5+5$ y por encima de $3-x/3$ y $(7x-17)/2$. El borde vertical está en $x=0$.
\begin{center}
\setlength{\unitlength}{0.9mm}
\begin{picture}(93,65)(-15,-8)
\qbezier(0,0)(37,0)(74,0)
\qbezier(0,0)(0,25)(0,50)
\qbezier(0,15)(0,20)(0,25)
\qbezier(0,25)(30,35)(60,45)
\qbezier(60,45)(48,27.5)(36,10)
\qbezier(36,10)(18,12.5)(0,15)
\put(75,0){$x$}\put(2,50){$y$}
\put(-13,14){\scriptsize $(0,3)$}\put(-13,25){\scriptsize $(0,5)$}
\put(36,4){\scriptsize $(3,2)$}\put(61,45){\scriptsize $(5,9)$}
\end{picture}
\end{center}
\textbf{b)} En $x=0$ quedan $(0,3)$ y $(0,5)$. La cota inferior ascendente se cruza con la otra inferior en $(3,2)$ y con la superior en $(5,9)$. Esos son los cuatro vértices.

\textbf{c)} En $(0,3),(0,5),(3,2),(5,9)$ los valores de $F$ son $3,1,10,7$. Por linealidad, el mínimo es $\boxed{1}$ sólo en $(0,5)$ y el máximo es $\boxed{10}$ sólo en $(3,2)$.
